\section{Interpretation and an experimental test}
\label{sec:interp}

\subsection{Parameter provenance and observation}

Every numerical rate in Section~\ref{sec:source} is assumed. Basal writing,
feedback, recruited erasure and division allocation specify an inherited-state
mechanism; they are not measured phenotype-switching rates. The
phenotype-to-death map is an explicit functional assumption, and the larger-site
sources of Section~\ref{sec:memory} are further mathematical realizations, not
more accurately calibrated loci. A physical larger domain might use a different
protected-state threshold, enzyme limitation, cooperative recruitment or site
dependence; those would be different sources and should be named as such.

Experiments would need separate division and death event records, calibrated
target engagement, a state observation model and \emph{joint} daughter
observations. Net growth $b-d$ alone is insufficient: the single-type sources
$(b,d)=(.1,.05)$ and $(.2,.15)$ share net growth $.05$ but have eventual
extinction probabilities $1/2$ and $3/4$. A binary reporter has emission rank at
most two and cannot satisfy full-column-rank identification of six arbitrary
states; this blocks that particular identification argument, not the possibility
of identifying a constrained low-dimensional rate family or a
certificate-relevant functional. No numerical range from another cell line or
treatment model is transplanted into this source.

\subsection{Conditional concentration and administered-amount bounds}
\label{sec:pkpd}

Suppose the target-engaging concentration $C(t)\ge0$ has a calibrated
saturating effect
\[
 v(C)=v_{\max}\frac{C}{K+C},\qquad K,v_{\max}>0 .
\]
Here $C$ and $K$ carry concentration units, $v_{\max}$ inverse time, and the
reaction exposure $B_{\mathrm{actual}}=\int_0^Tv(C(t))\,dt$ is dimensionless.
For finite concentration AUC $A=\int_0^TC(t)\,dt$, concavity and Jensen give
\[
 B_{\mathrm{actual}}\le\frac{v_{\max}TA}{KT+A}.
\]
If \eqref{eq:necessary} requires $B_{\mathrm{actual}}\ge B_{\mathrm{req}}>0$,
then
\begin{equation}\label{eq:auc}
 v_{\max}T>B_{\mathrm{req}},\qquad
 A\ge\frac{KTB_{\mathrm{req}}}{v_{\max}T-B_{\mathrm{req}}} .
\end{equation}
The strict inequality holds because finite $A$ makes $A/(KT+A)<1$;
multiplication by positive denominators proves the second. A concentration cap
adds $B_{\mathrm{req}}\le Tv(C_{\max})$.

For $\dot C=I/V_d-kC$ with $C(0)=0$, nonnegative infusion $I$, elimination
$k>0$ and distribution volume $V_d>0$, integration gives
$M=V_d(C(T)+kA)\ge kV_dA$, where $M=\int I$ is the administered amount, so
\begin{equation}\label{eq:mass}
 M\ \ge\ \frac{kV_dKTB_{\mathrm{req}}}{v_{\max}T-B_{\mathrm{req}}} .
\end{equation}
A nonzero initial concentration instead gives $M=V_d(C(T)-C(0)+kA)$ and
subtracts at most $V_dC(0)$ from the lower bound, truncated at zero. These are
necessary conditional bounds, not sufficient regimens; plasma concentration may
not equal the target-engaging concentration; and for a Hill exponent above one
global concavity fails, so one must use a valid concave majorant on the admitted
range rather than reusing this Jensen argument. The inequality prices the eraser
only, not the background killing retained in the example.

The useful content is qualitative and available before calibration: required
engagement close to $v_{\max}T$ forces the necessary AUC to increase sharply.
Section~\ref{sec:memory} sharpens the point. The amplitude floor is a constraint
on $v_{\max}$ alone, and no amount of AUC substitutes for it: if
$v_{\max}<v_c$, \eqref{eq:auc} is not merely hard to satisfy, it is irrelevant,
because Theorem~\ref{thm:amp} rules out eradication at any $A$ and any $T$.

\subsection{A falsifiable low-density comparison}

A prospective microcolony design can compare known founder counts $1,10,100$,
with $AA$-rich and $RR$-rich preparations, equal reaction-exposure early, late
and spread schedules, and separately declared eraser-only and
complete-withdrawal arms. Division and death event tracking should accompany
the molecular readout, with all founders retained in denominators. A fixed-time
viable-cell census and outgrowth over a declared observation window are
different events; undetected cells cannot be equated with extinct lineages
without a validated detection model.

The prediction is a one-sided requirement on exposure, not equality with a
fitted boundary. With a jointly valid parameter and preparation uncertainty
set, the lower bound is minimized and any policy upper bound maximized over the
set before comparing with data. For a simple fixed protocol and independent
replicates, a binomial confidence interval for the correctly observed event can
test that propagated prediction, with multiple-look error controlled in advance.
Survival significantly below a valid lower bound would challenge the adopted
source, actuator law or observation mapping; a data point above the lower bound,
or away from the sufficient upper curve, does not falsify it.

\subsection{Density, space and what the model is not}

The branching model is naturally a low-density persistence or microcolony
model. It is not automatically a model of a macroscopic tumour, and the
large-$n$ statements are mathematical scaling, not a claim that a
million-cell experiment remains density independent. Some extensions reuse the
barrier directly: for a population-dependent source it suffices to prove the
analogue of \eqref{eq:barrier} uniformly over the admitted population states and
environments, and a constructive upper bound survives if the weighted drift
remains uniformly negative. Reducing division is not automatically favourable
here, because division also redistributes marks and the sign of
$D_i(x,x)-x_i$ need not be uniform; density-dependent birth suppression
therefore cannot be inserted by an unsupported monotonicity argument. Spatial
migration or hidden carriers require explicit additional state variables and
generator terms, and persistent immigration destroys the absorbing-zero setup
entirely.
